\Needspace{9\baselineskip}
\begin{center}\textit{Lösung der Gleichung \(axx+byy+czz=0\).}\end{center}
\Article{294}
\textsc{Satz.} \textit{Seien \(a,b,c\) paarweise teilerfremde Zahlen,
von denen keine \(=0\) oder durch ein Quadrat teilbar ist. Die
Gleichung}
\[
axx+byy+czz=0\ldots(\Omega)
\]
\textit{besitzt keine ganzzahlige Lösung (außer \(x=y=z=0\), die wir
nicht berücksichtigen), sofern nicht \(-bc,-ac,-ab\) quadratische Reste
von \(a,b,c\) der Reihe nach sind und diese Zahlen nicht alle dasselbe Vorzeichen
haben. Sind diese vier Bedingungen aber erfüllt, so ist \((\Omega)\)
in ganzen Zahlen lösbar.}

\textit{Beweis.} Ist \((\Omega)\) überhaupt in ganzen Zahlen lösbar,
so auch durch Werte von \(x,y,z\), die keinen gemeinsamen Teiler
haben. Denn beliebige Werte, die der Gleichung \(\Omega\) genügen,
genügen ihr auch nach Division durch ihren größten gemeinsamen
Teiler. Nehmen wir nun \(app+bqq+crr=0\) an und seien \(p,q,r\) ohne gemeinsamen
Teiler, so sind sie auch paarweise teilerfremd. Hätten nämlich
\(q,r\) einen gemeinsamen Teiler \(\mu\), so wäre dieser zu \(p\)
teilerfremd, aber \(\mu\mu\) würde \(app\) und somit auch \(a\)
teilen, entgegen der Voraussetzung. Ebenso sind \(p,r;\ p,q\) paarweise
teilerfremd. Daher wird \(-app\) durch die binäre Form \(byy+czz\)
\SourceLeaf{355}
dargestellt, indem man \(y,z\) die teilerfremden Werte \(q,r\)
gibt. Ihre Determinante \(-bc\) ist somit quadratischer Rest
von \(app\) und daher auch von \(a\) (Art.\ 154). Ebenso ist
\(-acRb,-abRc\). Dass aber \((\Omega)\) keine Lösung haben kann, wenn \(a,b,c\)
dasselbe Vorzeichen haben, ist so offensichtlich, dass es keiner
Erklärung bedarf.

Den Beweis der Umkehrung, die den zweiten Teil des Satzes bildet,
richten wir so ein: \textit{Zuerst} zeigen wir, wie man eine zu
\(\left(\begin{smallmatrix}a,b,c\\0,0,0\end{smallmatrix}\right)\ldots f\) äquivalente ternäre Form findet, deren Koeffizienten
\(2,3,4\) durch \(abc\) teilbar sind. Daraus leiten wir \textit{zweitens}
eine Lösung der Gleichung \((\Omega)\) ab.

I.\ Man bestimme drei ganze Zahlen \(A,B,C\) ohne gemeinsamen Teiler
so, dass \(A\) zu \(b\) und \(c\), \(B\) zu \(a\) und
\(c\), \(C\) zu \(a\) und \(b\) teilerfremd ist, während
\(aAA+bBB+cCC\) durch \(abc\) teilbar ist. Dies erreicht man wie folgt.
Seien \(\mathfrak A,\mathfrak B,\mathfrak C\) Werte der Ausdrücke \(\surd-bc\pmod a\), \(\surd-ac\pmod b\), \(\surd-ab\pmod c\),
die notwendig zu \(a,b,c\) der Reihe nach teilerfremd sind. Man
wähle drei ganze Zahlen \(\mathfrak a,\mathfrak b,\mathfrak c\) ganz beliebig, nur so, dass sie
der Reihe nach zu \(a,b,c\) teilerfremd sind (\textit{zum Beispiel}
alle \(=1\)), und bestimme \(A,B,C\) so, dass
\[
\begin{aligned}
A&\equiv\mathfrak b c\pmod b&&\text{und }\equiv\mathfrak c\mathfrak C\pmod c\\
B&\equiv\mathfrak c a\pmod c&&\text{und }\equiv\mathfrak a\mathfrak A\pmod a\\
C&\equiv\mathfrak a b\pmod a&&\text{und }\equiv\mathfrak b\mathfrak B\pmod b
\end{aligned}
\]
Dann ist
\[
\begin{aligned}
aAA+bBB+cCC&\equiv\mathfrak a\mathfrak a(b\mathfrak A\mathfrak A+cbb)\\
&\equiv\mathfrak a\mathfrak a(b\mathfrak A\mathfrak A-\mathfrak A\mathfrak A b)\equiv0\pmod a
\end{aligned}
\]
also durch \(a\) teilbar, ebenso durch \(b,c\) und somit auch
durch \(abc\). Außerdem ist offenbar \(A\) notwendig zu \(b\)
und \(c\), \(B\) zu \(a\) und \(c\), \(C\) zu \(a\)
und \(b\) teilerfremd. Haben diese Werte von \(A,B,C\) aber einen
(größten) gemeinsamen Teiler \(\mu\), so ist dieser offenbar
zu \(a,b,c\) und somit zu \(abc\) teilerfremd. Durch Division
jener Werte durch \(\mu\) erhalten wir daher neue ohne gemeinsamen
Teiler, die für \(aAA+bBB+cCC\) weiterhin einen durch \(abc\) teilbaren
Wert ergeben und somit alle Bedingungen erfüllen.

II.\ Sind die Zahlen \(A,B,C\) so bestimmt, so haben auch \(Aa,Bb,Cc\)
keinen gemeinsamen Teiler. Hätten sie nämlich einen gemeinsamen
Teiler \(\mu\), so wäre dieser notwendig zu \(a\) teilerfremd
(da dieses sowohl zu \(Bb\) als auch zu \(Cc\) teilerfremd ist)
und ebenso zu
\SourceLeaf{356}
\(b\) und \(c\). Daher müsste \(\mu\) auch \(A,B,C\) teilen,
entgegen der Voraussetzung. Man kann also ganze Zahlen \(\alpha,\beta,\gamma\)
finden, sodass \(\alpha Aa+\beta Bb+\gamma Cc=1\) ist. Ferner bestimme man sechs ganze
Zahlen \(\alpha',\beta',\gamma',\alpha'',\beta'',\gamma''\) so, dass
\[
\beta'\gamma''-\gamma'\beta''=Aa,\quad
\gamma'\alpha''-\alpha'\gamma''=Bb,\quad
\alpha'\beta''-\beta'\alpha''=Cc
\]
Nun gehe \(f\) durch die Substitution
\[
\begin{array}{ccc}
\alpha,&\alpha',&\alpha''\\
\beta,&\beta',&\beta''\\
\gamma,&\gamma',&\gamma''
\end{array}
\]
in \(\left(\begin{smallmatrix}m,m',m''\\n,n',n''\end{smallmatrix}\right)=g\) über (die zu \(f\) äquivalent ist). Ich behaupte,
dass \(m',m'',n\) durch \(abc\) teilbar sind. Setzt man nämlich
\[
\begin{aligned}
\beta''\gamma-\gamma''\beta&=A',&\gamma''\alpha-\alpha''\gamma&=B',&\alpha''\beta-\beta''\alpha&=C'\\
\beta\gamma'-\gamma\beta'&=A'',&\gamma\alpha'-\alpha\gamma'&=B'',&\alpha\beta'-\beta\alpha'&=C''
\end{aligned}
\]
so ist
\[
\begin{aligned}
\alpha'&=B''Cc-C''Bb,&\beta'&=C''Aa-A''Cc,&\gamma'&=A''Bb-B''Aa\\
\alpha''&=C'Bb-B'Cc,&\beta''&=A'Cc-C'Aa,&\gamma''&=B'Aa-A'Bb
\end{aligned}
\]
Durch Einsetzen dieser Werte in die Gleichungen
\[
\begin{aligned}
m'&=a\alpha'\alpha'+b\beta'\beta'+c\gamma'\gamma'\\
m''&=a\alpha''\alpha''+b\beta''\beta''+c\gamma''\gamma''\\
n&=a\alpha'\alpha''+b\beta'\beta''+c\gamma'\gamma''
\end{aligned}
\]
erhält man modulo \(a\)
\[
\begin{aligned}
m'&\equiv bcA''A''(BBb+CCc)\equiv0\\
m''&\equiv bcA'A'(BBb+CCc)\equiv0\\
n&\equiv bcA'A''(BBb+CCc)\equiv0
\end{aligned}
\]
also sind \(m',m'',n\) durch \(a\) teilbar. Ebenso findet man,
dass dieselben Zahlen durch \(b,c\) und somit auch durch
\(abc\) teilbar sind, wie behauptet.
\SourceLeaf{357}
\par
III.\ Für die Determinante der Formen \(f,g\) schreiben wir
der Übersicht wegen \(-abc=d\) und setzen
\[
md=M,\quad m'=M'd,\quad m''=M''d,\quad n=Nd,\quad n'=N',\quad n''=N''
\]
Offenbar geht \(f\) durch die Substitution \((S)\)
\[
\begin{array}{ccc}
\alpha d,&\alpha',&\alpha''\\
\beta d,&\beta',&\beta''\\
\gamma d,&\gamma',&\gamma''
\end{array}
\]
in die ternäre Form \(\left(\begin{smallmatrix}Md,M'd,M''d\\Nd,N'd,N''d\end{smallmatrix}\right)=g'\) mit der Determinante \(d^3\)
über, die daher unter \(f\) enthalten ist. Nun behaupte ich,
dass diese Form \(g'\) notwendig zu \(\left(\begin{smallmatrix}d,0,0\\d,0,0\end{smallmatrix}\right)=g''\) äquivalent ist.
Denn \(\left(\begin{smallmatrix}M,M',M''\\N,N',N''\end{smallmatrix}\right)=g'''\) ist offenbar eine ternäre Form mit der Determinante
\(1\). Da außerdem nach Voraussetzung \(a,b,c\) nicht alle
dasselbe Vorzeichen haben, ist \(f\) eine indefinite Form.
Daraus folgt leicht, dass auch \(g'\) und \(g'''\) indefinit
sein müssen. Daher ist \(g'''\) zur Form \(\left(\begin{smallmatrix}1,0,0\\1,0,0\end{smallmatrix}\right)\) äquivalent
(Art.\ 277), und man kann eine Transformation \((S')\) der
ersten in die zweite finden. Offenbar geht aber durch \((S')\)
die Form \(g'\) in \(g''\) über. Somit ist auch \(g''\)
unter \(f\) enthalten, und aus der Verbindung der Substitutionen
\((S),(S')\) folgt eine Transformation der Form \(f\) in \(g''\).
Ist diese
\[
\begin{array}{ccc}
\delta,&\delta',&\delta''\\
\varepsilon,&\varepsilon',&\varepsilon''\\
\zeta,&\zeta',&\zeta''
\end{array}
\]
so erhält man offenbar zwei Lösungen der Gleichung \((\Omega)\),
nämlich \(x=\delta',y=\varepsilon',z=\zeta'\) und \(x=\delta'',y=\varepsilon'',z=\zeta''\). Zugleich ist klar, dass in keiner
der beiden Wertgruppen alle Werte \(=0\) werden können, da
notwendig
\[
\delta\varepsilon'\zeta''+\delta'\varepsilon''\zeta+\delta''\varepsilon\zeta'-\delta\varepsilon''\zeta'-\delta'\varepsilon\zeta''-\delta''\varepsilon'\zeta=d.\quad\text{Wie verlangt.}
\]

\emph{Beispiel.} Die gegebene Gleichung sei \(7xx-15yy+23zz=0\). Sie ist
lösbar, weil \(345\,R\,7,-161\,R\,15,105\,R\,23\). Hier sind die Werte von \(\mathfrak A,\mathfrak B,\mathfrak C\) gleich
\(3,7,6\); setzt man \(\mathfrak a=\mathfrak b=\mathfrak c=1\), so findet man \(A=98,B=-39,C=-8\). Daraus
erhält man die Substitution
\[
\left\{\begin{array}{rrr}3,&5,&22\\-1,&2,&-28\\8,&25,&-7\end{array}\right\}
\]
durch die \(f\) in \(\left(\begin{smallmatrix}1520,14490,-7245\\-2415,-1246,4735\end{smallmatrix}\right)=g\) übergeht. Daher ist
\[
(S)=\left\{\begin{array}{rrr}7245,&5,&22\\-2415,&2,&-28\\19320,&25,&-7\end{array}\right\},\quad
g'''=\left(\begin{smallmatrix}3670800,6,-3\\-1,-1246,4735\end{smallmatrix}\right).
\]
\SourceLeaf{358}
Man findet, dass die Form \(g'''\) in \(\left(\begin{smallmatrix}1,0,0\\1,0,0\end{smallmatrix}\right)\) durch die Substitution
\[
\left\{\begin{array}{rrr}3,&5,&1\\-2440,&-4066,&-813\\-433,&-722,&-144\end{array}\right\}\ldots(S')
\]
übergeht. Verbindet man diese mit \((S)\), so
ergibt sich \(\left\{\begin{smallmatrix}9,11.12\\-1,9,-9\\-9,4,3\end{smallmatrix}\right\}\), durch die \(f\) in \(g''\) übergeht.
Wir haben also zwei Lösungen der gegebenen Gleichung, \(x=11,y=9,z=4\)
und \(x=12,y=-9,z=3\). Die zweite wird einfacher, wenn man die Werte
durch den gemeinsamen Teiler \(3\) dividiert; daraus folgt
\(x=4,y=-3,z=1\).

\Article{295}
Den zweiten Teil des Satzes im vorigen Artikel kann man auch
wie folgt beweisen. Man suche eine ganze Zahl \(h\) so,
dass \(ah\equiv\mathfrak C\pmod c\) ist (die Zeichen \(\mathfrak A,\mathfrak B,\mathfrak C\) verwenden wir im selben
Sinn wie im vorigen Artikel), und setze \(ahh+b=ci\). Dann sieht
man leicht, dass \(i\) ganzzahlig und die Zahl \(-ab\) die
Determinante der binären Form \((ac,ah,i)\ldots\varphi\) ist. Diese Form ist
sicher nicht positiv (denn da nach Voraussetzung \(a,b,c\)
nicht alle dasselbe Vorzeichen haben, können \(ab\) und
\(ac\) nicht zugleich positiv sein). Ferner hat sie die
charakteristische Zahl \(-1\), was wir wie folgt synthetisch
beweisen: Man bestimme ganze Zahlen \(e,e'\) so, dass
\[
\begin{aligned}
e&\equiv0\pmod a\text{ und }\equiv\mathfrak B\pmod b;\\
ce'&\equiv\mathfrak A\pmod a\text{ und }\equiv h\mathfrak B\pmod b
\end{aligned}
\]
Dann ist \((e,e')\) ein Wert des Ausdrucks \(\surd-(ac,ah,i)\pmod{-ab}\). Denn modulo
\(a\) ist
\[
\begin{aligned}
ee&\equiv0\equiv-ac,\quad ee'\equiv0\equiv-ah\\
cce'e'&\equiv\mathfrak A\mathfrak A\equiv-bc\equiv-cci\quad\text{also }e'e'\equiv-i
\end{aligned}
\]
modulo \(b\) dagegen
\[
\begin{aligned}
ee&\equiv\mathfrak B\mathfrak B\equiv-ac,\quad cee'\equiv h\mathfrak B\mathfrak B\equiv-ach\quad\text{also }ee'\equiv-ah\\
cce'e'&\equiv hh\mathfrak B\mathfrak B\equiv-achh\equiv-cci\quad\text{also }e'e'\equiv-i
\end{aligned}
\]
Dieselben drei Kongruenzen, die für beide Moduln \(a,b\)
gelten, gelten aber auch modulo \(ab\). Aus der Theorie
der ternären Formen folgt leicht, dass \(\varphi\) durch die
Form \(\left(\begin{smallmatrix}-1,0,0\\1,0,0\end{smallmatrix}\right)\) darstellbar ist. Es sei also
\[
actt+2ahtu+iuu=-(\alpha t+\beta u)^2+2(\gamma t+\delta u)(\varepsilon t+\zeta u)
\]
\SourceLeaf{359}
Durch Multiplikation mit \(c\) erhält man
\[
a(ct+hu)^2+buu=-c(\alpha t+\beta u)^2+2c(\gamma t+\delta u)(\varepsilon t+\zeta u)
\]
Gibt man nun \(t,u\) bestimmte Werte, sodass entweder
\(\gamma t+\delta u\) oder \(\varepsilon t+\zeta u\) \(=0\) wird, so hat man
eine Lösung der Gleichung \(\Omega\). Dieser genügen daher
sowohl
\[
x=\delta c-\gamma h,\quad y=\gamma,\quad z=\alpha\delta-\beta\gamma
\]
als auch
\[
x=\zeta c-\varepsilon h,\quad y=\varepsilon,\quad z=\alpha\zeta-\beta\varepsilon
\]
Dabei ist klar, dass weder die ersten noch die zweiten
Werte sämtlich \(=0\) werden können. Denn wäre \(\delta c-\gamma h=0,\gamma=0\),
so wäre auch \(\delta=0\) und \(\varphi=-(\alpha t+\beta u)^2\), woraus \(ab=0\) folgte,
entgegen der Voraussetzung; ebenso bei den anderen Werten.
--- In unserem Beispiel finden wir als Form \(\varphi\) die
Form \((161,-63,24)\), als Wert des Ausdrucks \(\surd-\varphi\pmod{105}=(7,-51)\) und als
Darstellung der Form \(\varphi\) durch \(\left(\begin{smallmatrix}-1,0,0\\1,0,0\end{smallmatrix}\right)\) die folgende:
\[
\varphi=-(13t-4u)^2+2(11t-4u)(15t-5u)
\]
Daraus folgen die Lösungen \(x=7,y=11,z=-8\); \(x=20,y=15,z=-5\), oder nach
Division durch \(5\) und Vernachlässigung des Vorzeichens
von \(z\) die Lösung \(x=4,y=3,z=1\).

Von diesen beiden Methoden zur Lösung der Gleichung \(\Omega\)
hat die zweite den Vorteil, dass man meist mit kleineren
Zahlen auskommt. Die erste, die sich auch durch verschiedene
hier übergangene Kunstgriffe verkürzen lässt, erscheint aber
eleganter, besonders weil die Zahlen \(a,b,c\) ganz gleich
behandelt werden und sich die Rechnung durch keine Vertauschung
dieser Zahlen ändert. Anders bei der zweiten Methode: Hier
ergibt sich meist die bequemste Rechnung, wenn man für
\(a\) die kleinste und für \(c\) die größte der drei
gegebenen Zahlen nimmt, wie wir es in unserem Beispiel
getan haben.

\begin{center}\itshape Über die Methode, mit der der berühmte Le Gendre den Fundamentalsatz behandelte.\end{center}
\Article{296}
Den eleganten Satz der vorigen Artikel hat zuerst der
berühmte Le Gendre gefunden, \emph{Hist. de l'Ac. de Paris}
1784, S.\ 507, und mit einem schönen, von unseren beiden
ganz verschiedenen Beweis versehen. Zugleich versuchte
dieser ausgezeichnete Geometer dort, daraus einen Beweis
von Sätzen abzuleiten, die mit dem Fundamentalsatz
\SourceLeaf{360}
des vorigen Abschnitts übereinstimmen. Dass uns dieser
Beweis für diesen Zweck ungeeignet erscheint, haben wir
schon oben in Art.\ 151 erklärt. Hier ist daher der Ort,
diesen an sich sehr eleganten Beweis kurz darzustellen
und die Gründe für unser Urteil hinzuzufügen. Wir schicken
die folgende Bemerkung voraus: \emph{Sind die Zahlen
\(a,b,c\) sämtlich \(\equiv1\pmod4\), so kann die Gleichung \(axx+byy+czz=0\ldots(\Omega)\)
nicht lösbar sein.} Denn man sieht sehr leicht, dass der
Wert von \(axx+byy+czz\) in diesem Fall notwendig entweder
\(\equiv1\), \(\equiv2\) oder \(\equiv3\pmod4\) ist, sofern nicht alle
\(x,y,z\) zugleich gerade gewählt werden. Wäre also
\(\Omega\) lösbar, so nur durch gerade Werte von \(x,y,z\),
\emph{was ein Widerspruch ist}, da beliebige Werte, die der
Gleichung \(\Omega\) genügen, ihr auch nach Division durch
ihren größten gemeinsamen Teiler genügen, wobei notwendig
mindestens ein ungerader Wert herauskommen muss. Die
verschiedenen Fälle des zu beweisenden Satzes lassen
sich nun auf die folgenden Punkte zurückführen.

I.\ Seien \(p,q\) Primzahlen der Form \(4n+3\) (positiv
und ungleich). Dann kann nicht zugleich \(p\,R\,q,q\,R\,p\) sein.
Wäre dies nämlich möglich, so wären mit \(1=a,-p=b,-q=c\) offenbar
alle Bedingungen für die Lösbarkeit der Gleichung \(axx+byy+czz=0\)
erfüllt (Art.\ 294). Nach der vorigen Bemerkung hat sie
aber keine Lösung; die Annahme kann also nicht bestehen.
Daraus folgt sofort Satz 7 in Art.\ 131.

II.\ Ist \(p\) eine Primzahl der Form \(4n+1\) und
\(q\) eine Primzahl der Form \(4n+3\), so kann nicht
zugleich \(q\,R\,p,p\,N\,q\) sein. Sonst wäre nämlich \(-p\,R\,q\) und
die Gleichung \(xx+pyy-qzz=0\) lösbar, die nach der vorigen
Bemerkung keine Lösung zulässt. Daraus folgen die Fälle
4 und 5 in Art.\ 131.

III.\ Sind \(p,q\) Primzahlen der Form \(4n+1\), so
kann nicht zugleich \(p\,R\,q,q\,N\,p\) sein. Man nehme eine andere
Primzahl \(r\) der Form \(4n+3\), die Rest von \(q\)
ist und von der \(p\) Nichtrest ist. Nach den gerade
bewiesenen Fällen (II) ist dann \(q\,R\,r,r\,N\,p\). Wäre also
\(p\,R\,q,q\,N\,p\), so wäre \(qr\,R\,p,pr\,R\,q,pq\,N\,r\) und daher \(-pq\,R\,r\). Somit
wäre die Gleichung \(pxx+qyy-rzz=0\) lösbar, entgegen der vorigen
Bemerkung; die Annahme kann also nicht bestehen. Daraus
folgen die Fälle 1 und 2 in Art.\ 131.

Dieser Fall lässt sich auf folgende Weise übersichtlicher
behandeln. Sei \(r\) eine Primzahl der Form \(4n+3\),
von der \(p\) Nichtrest ist. Dann ist auch \(r\,N\,p\),
also (unter der Annahme \(p\,R\,q,q\,N\,p\)) \(qr\,R\,p\). Ferner ist
\(-p\,R\,q,-p\,R\,r\) und somit auch \(-p\,R\,qr\); daher wäre die Gleichung
\(xx+pyy-qrzz=0\) lösbar, entgegen der vorigen Bemerkung. Daraus
usw.

IV.\ Ist \(p\) eine Primzahl der Form \(4n+1\) und
\(q\) eine Primzahl der Form \(4n+3\),
\SourceLeaf{361}
so kann nicht zugleich \(p\,R\,q,q\,N\,p\) sein. Man nehme eine
Hilfsprimzahl \(r\) der Form \(4n+1\), die Nichtrest
von beiden \(p,q\) ist. Dann ist (nach II) \(q\,N\,r\)
und (nach III) \(p\,N\,r\), also \(pq\,R\,r\). Wäre daher
\(p\,R\,q,q\,N\,p\), so wäre auch \(pr\,N\,q,-pr\,R\,q,qr\,R\,p\) und die Gleichung
\(pxx-qyy+rzz=0\) lösbar, \emph{was ein Widerspruch ist}. --- Daraus
folgen die Fälle 3 und 6 in Art.\ 131.

V.\ Seien \(p,q\) Primzahlen der Form \(4n+3\). Dann
kann nicht zugleich \(p\,N\,q,q\,N\,p\) sein. Angenommen, dies wäre
möglich, und man nehme eine Hilfsprimzahl \(r\) der
Form \(4n+1\), die Nichtrest von beiden \(p,q\) ist.
Dann wäre \(qr\,R\,p,pr\,R\,q\), ferner (nach II) \(p\,N\,r,q\,N\,r\), also
\(pq\,R\,r\) und \(-pq\,R\,r\). Somit wäre die Gleichung \(-pxx-qyy+rzz=0\)
möglich, entgegen der vorigen Bemerkung. Daraus folgt
Fall 8 in Art.\ 131.
